\chapter{Полевое суперсвязностное завершение цепи когерентности}
\label{ch:v7-edge-coherence-field-space-superconnection-gate}

\section{Разделение физических и вычислительных носителей}

Бимодульный гейт запретил отождествлять цепь $1\to6\to3$ с новыми
фермионными вершинами. Это не запрещает использовать ту же цепь как
ассоциированный градуированный пучок над пространством полей. Пусть
\begin{equation}
 V=\mathbb C^2_{\rm copy},\qquad
 W=W_e\oplus W_X\oplus W_Y\simeq\mathbb C^3,
 \qquad B\in\operatorname{Hom}(W,V).
 \label{eq:v7-field-superconnection-arrow-bundle}
\end{equation}
Разложение $W$ сохраняется явно: каналы $e_R,X_R$ изотипичны, тогда как
$Y_R$ имеет другой бимодульный тип. Поэтому полный $U(3)$ на $W$ является
лишь случайной симметрией абстрактной нормы и не объявляется физической
калибровочной группой.

Введём ассоциированный градуированный носитель
\begin{equation}
 \mathcal E^0=\underline{\mathbb C},\qquad
 \mathcal E^1=\operatorname{Hom}(W,V),\qquad
 \mathcal E^2=\operatorname{Hom}(\Lambda^2W,\Lambda^2V).
 \label{eq:v7-field-superconnection-carrier}
\end{equation}
Его волокна имеют размерности $1,6,3$, но секции $\mathcal E^k$ не
интерпретируются как фермионные поля. Это вычислительные волокна,
функториально построенные из пространства амплитуд разрешённых стрелок.

\section{Ориентированная суперсвязность и её кривизна}

Определим нечётный оператор положительной степени
\begin{equation}
 d_B:\mathcal E^0\oplus\mathcal E^1\oplus\mathcal E^2
 \longrightarrow
 \mathcal E^0\oplus\mathcal E^1\oplus\mathcal E^2,
 \qquad
 d_B|_{\mathcal E^0}=A_B,\quad
 d_B|_{\mathcal E^1}=C_B,\quad
 d_B|_{\mathcal E^2}=0,
 \label{eq:v7-field-superconnection-oriented-differential}
\end{equation}
где $A_B(1)=B$ и $C_B=\frac12d(\Lambda^2)_B$. Тогда
\begin{equation}
 d_B^2|_{\mathcal E^0}=C_BA_B=\Lambda^2B,qquad
 d_B^2|_{\mathcal E^1\oplus\mathcal E^2}=0.
 \label{eq:v7-field-superconnection-curvature}
\end{equation}
Следовательно,
\begin{equation}
 \|d_B^2\|^2=\|\Lambda^2B\|_F^2
 =\det(BB^\dagger).
 \label{eq:v7-field-superconnection-curvature-norm}
\end{equation}
Вне вакуума это кривой комплекс: квадрат дифференциала является
алгебраической кривизной. На ранговой страте
\begin{equation}
 \rank B\leq1
 \quad\Longleftrightarrow\quad
 d_B^2=0
 \label{eq:v7-field-superconnection-integrable-locus}
\end{equation}
он становится настоящим комплексом. Это не означает исчезновения полной
эрмитовой кривизны: условие относится именно к ориентированной части
$d_B$, а не к $\mathcal D_B=d_B+d_B^\dagger$.

Пусть $\nabla^V$ и $\nabla^W$ --- уже имеющиеся связности конечных
endpoint-пучков. На двух составных волокнах используются только
индуцированные связности
\begin{align}
 \nabla^{\mathcal E^1}X
 &=\nabla^V X-X\nabla^W,\\
 \nabla^{\mathcal E^2}Y
 &=\nabla^{\Lambda^2V}Y-Y\nabla^{\Lambda^2W}.
 \label{eq:v7-field-superconnection-induced-connections}
\end{align}
Таким образом, суперсвязность $\mathbb A=\nabla^{\mathcal E}+d_B$ не
требует независимой калибровочной связности на узлах размерностей $6$ и
$3$.

\section{Фактическая блочная ковариантность}

Максимальная группа замен базиса, сохраняющая бимодульную типизацию
каналов, имеет вид
\begin{equation}
 G_{\rm blk}=U(2)_{\rm copy}\times U(2)_{eX}\times U(1)_Y
 \subset U(V)\times U(W).
 \label{eq:v7-field-superconnection-block-group}
\end{equation}
Для $g\in U(V)$ и блочно-диагонального $h\in U(W)$ поле преобразуется как
\begin{equation}
 B\longmapsto gBh^\dagger.
 \label{eq:v7-field-superconnection-field-action}
\end{equation}
Функториальность внешней степени даёт
\begin{align}
 \Lambda^2(gBh^\dagger)
 &=\Lambda^2g\,(\Lambda^2B)\,(\Lambda^2h)^\dagger,\\
 C_{gBh^\dagger}(gXh^\dagger)
 &=\Lambda^2g\,C_B(X)\,(\Lambda^2h)^\dagger.
 \label{eq:v7-field-superconnection-covariance}
\end{align}
Поэтому ориентированная и эрмитова части ковариантны относительно
$G_{\rm blk}$ и, следовательно, относительно действующей физической
калибровочной подгруппы. Для результата не требуется смешивать $Y_R$ с
$e_R,X_R$ фиктивным физическим $U(3)$.

\section{Эрмитова часть, потенциал и кинетическая метрика}

Эрмитова нечётная часть
\begin{equation}
 \mathcal D_B=d_B+d_B^\dagger
 \label{eq:v7-field-superconnection-hermitian-part}
\end{equation}
совпадает с оператором спектрального родителя. Поэтому без добавления
новых полей сохраняются тождества
\begin{align}
 \Tr\mathcal D_B^2&=3\Tr(BB^\dagger),\\
 \Tr\mathcal D_B^4&=\frac94\bigl(\Tr BB^\dagger\bigr)^2
 +\frac{15}{4}\det(BB^\dagger).
 \label{eq:v7-field-superconnection-finite-traces}
\end{align}
При пространственно-временной зависимости $B(x)$ главная квадратичная
форма производных фиксируется той же конечной следовой метрикой:
\begin{equation}
 \boxed{
 \Tr\bigl(\delta\mathcal D_B\,\delta\mathcal D_B\bigr)
 =3\Tr(\delta B\,\delta B^\dagger).}
 \label{eq:v7-field-superconnection-kinetic-metric}
\end{equation}
На двенадцати вещественных компонентах её матрица равна $3I_{12}$.
Следовательно, полевой подъём не создаёт отрицательных кинетических мод.
Абсолютный коэффициент пространственно-временного кинетического члена
зависит от теплового разложения полного произведения и здесь не считается
выведенным.

Real-завершение строится сопряжённым ассоциированным пучком и антилинейным
обменом. Полуслед Real-пары возвращает ровно исходные второй и четвёртый
моменты, поэтому скрытый множитель два не возникает.

\begin{theorem}[Допуск полевого суперсвязностного носителя]
Цепь $1\to6\to3$ канонически реализуется как ассоциированный
градуированный пучок пространства стрелок, не добавляя физических
фермионных вершин и независимых калибровочных связностей. Ориентированная
кривизна равна $\Lambda^2B$ и исчезает точно на страте ранга не выше
одного. Эрмитова часть сохраняет единый спектральный потенциал, а её
следовая кинетическая метрика положительна и равна $3I_{12}$. Конструкция
ковариантна относительно фактической блочной группы без постулирования
физического $U(3)$ каналов.
\end{theorem}

Этот результат закрывает вопрос о математическом носителе положительного
механизма, но ещё не завершает физический селектор. Поле $B$ занимает
шестикромочный прямоугольник полного строгого графа. Внутри него находится
одно старое и пять новых рёбер; только два из этих пяти принадлежат целевому
ремонту, а три являются нежелательными. Ещё шесть новых рёбер лежат вне
$B$. Поэтому прямоугольник когерентности не совпадает с требуемым меню из
шести новых блоков.

\begin{nogocriterion}[Запрет скрытого проектора на шесть рёбер]
Нельзя считать полевой суперсвязностный допуск полным выбором физического
графа, пока явно не сопоставлены его ранга-один вакуум и точное целевое
меню рёбер, а все одиннадцать новых разрешённых блоков не включены в один
тест. Если потребуется заранее занулить конкурирующие рёбра или вставить
проектор на выбранный рисунок, успех останется условным.
\end{nogocriterion}

\begin{protocol}
\begin{enumerate}
 \item физические фермионные вершины: \textbf{не добавлены};
 \item вычислительный носитель $1,6,3$: \textbf{ассоциирован функториально};
 \item ориентированная кривизна: \textbf{$d_B^2=\Lambda^2B$};
 \item интегрируемая страта: \textbf{$\rank B\leq1$};
 \item независимые калибровочные связности: \textbf{не добавлены};
 \item фактическая блочная ковариантность: \textbf{подтверждена};
 \item физический $U(3)$ каналов: \textbf{не постулируется};
 \item спектральный потенциал: \textbf{сохранён};
 \item кинетическая метрика: \textbf{$3I_{12}$, положительна};
 \item Real-полуслед: \textbf{сохраняет коэффициенты};
 \item абсолютная тепловая нормировка: \textbf{открыта};
 \item конкуренция всех одиннадцати рёбер: \textbf{открыта};
 \item итог: \textbf{положительный вспомогательный полевой носитель};
 \item следующий гейт: \textbf{полное одиннадцатирёберное вложение без
 ручного проектора}.
\end{enumerate}
\end{protocol}

Машинный сертификат сохранён в
\path{s2t_v7_edge_coherence_field_space_superconnection_gate_results.json}.

\begin{thebibliography}{9}
\bibitem{v7-field-superconnection-quillen}
D.~Quillen,
\emph{Superconnections and the Chern character},
Topology 24 (1985), 89--95.
\bibitem{v7-field-superconnection-exterior}
G.~Roepstorff,
\emph{Superconnections and Matter}, arXiv:hep-th/9801040.
\bibitem{v7-field-superconnection-correspondence}
B.~Mesland,
\emph{Unbounded bivariant K-theory and correspondences in noncommutative
geometry}, arXiv:0904.4383.
\bibitem{v7-field-superconnection-gauge-product}
S.~Brain, B.~Mesland, W.~D.~van Suijlekom,
\emph{Gauge theory for spectral triples and the unbounded Kasparov product},
Journal of Noncommutative Geometry 10 (2016), 135--206;
arXiv:1306.1951.
\end{thebibliography}